\section*{Schiffmann--Vasserot at chain level: $\PhiFA_3(\Perf(\C^3)) = \CoHA(\C^3) = Y^+(\widehat{\fgl}_1)$ as $E_3$ factorisation algebra on $\Hilb^\bullet(\C^3)$}
\label{sec:sv-chain-level-e3-factorization-cfg2026}

Fix the following data. Let $V = \C^3$ with its standard torus action
by $T = (\C^\times)^3$. Let $\epsilon_1, \epsilon_2, \epsilon_3 \in
H^*_T(\pt; \Q) = \Q[\epsilon_1, \epsilon_2, \epsilon_3]$ be the
equivariant weights of the three coordinate directions, with the
Calabi--Yau constraint
\[
 \epsilon_1 + \epsilon_2 + \epsilon_3 = 0
\]
cutting out the two-parameter sub-$T$ preserving the holomorphic
$3$-form $dx \wedge dy \wedge dz$. Write $\hbar = -\epsilon_1
\epsilon_2 \epsilon_3$ (a cubic invariant of the full $(\epsilon_1,
\epsilon_2, \epsilon_3)$-algebra) and $\kappa = \epsilon_1 + \epsilon_2
+ \epsilon_3$ (which vanishes on the CY locus; in the full
Schiffmann--Vasserot parameter space $\kappa$ is a free deformation
parameter, and the CY specialisation is $\kappa = 0$).

\subsection*{The Hilbert scheme of points on $\C^3$ and its $E_3$ factorisation structure}

\begin{definition}[Hilbert scheme $\Hilb^n(\C^3)$]
\label{def:hilb-c3-n}
\ClaimStatusDefinitional
For $n \geq 0$, $\Hilb^n(\C^3)$ is the moduli scheme of ideals
$I \subset \C[x, y, z]$ with $\dim_\C \C[x, y, z] / I = n$; equivalently,
length-$n$ subschemes of $\C^3$. It is a quasi-projective $\C$-variety;
at $n = 0$ it is a point; at $n = 1$ it is $\C^3$ itself; for $n \geq 2$
it is singular, carrying an obstructed deformation theory
(Nakajima~\textit{Lectures on Hilbert schemes of points on surfaces},
American Mathematical Society $1999$; Nakajima--Yoshioka~$2003$,
\textit{Invent. Math.} $155$; Maulik--Nekrasov--Okounkov--Pandharipande
$\mathrm{I}$~$2006$, arXiv:math/$0312059$, \S$1.6$). Write
$\Hilb^\bullet(\C^3) = \bigsqcup_{n \geq 0} \Hilb^n(\C^3)$ for the
disjoint-union moduli of all ideal lengths.
\end{definition}

\begin{remark}[Torus action]
\label{rem:hilb-c3-torus-action}
The torus $T = (\C^\times)^3$ acts on $\C^3$ by scaling the three
coordinates. The induced action on $\Hilb^n(\C^3)$ has isolated fixed
points indexed by monomial ideals of colength $n$, equivalently by
three-dimensional partitions (plane partitions) $\pi$ of size $n$
(Okounkov--Reshetikhin $2003$, \textit{Duke Math. J.} $119$;
Nakajima~\textit{Lectures}, \S$9.1$ for the two-dimensional analogue
on $\C^2$). At a monomial-ideal fixed point $I_\pi$, the tangent space
$T_{I_\pi} \Hilb^n(\C^3)$ is a direct sum of $T$-weight spaces with
weights
\[
 \bigl\{\epsilon_1 l(s) - \epsilon_2 (a(s) + 1),\;
 -\epsilon_1 (l(s) + 1) + \epsilon_2 a(s),\;
 \epsilon_3 \cdot (\text{height shifts})\bigr\}_{s \in \pi}
\]
with $l(s), a(s)$ the leg and arm of a box in the two-dimensional
slice and the third factor accounting for the $z$-direction. The
explicit character is computed in MNOP~$\mathrm{I}$ Theorem~$1$. The
equivariant cohomology $H^*_T(\Hilb^n(\C^3); \Q)$ has a basis $\{[I_\pi]\}$
indexed by plane partitions of size $n$, with localisation formula
\[
 \alpha \;=\; \sum_\pi \frac{\iota_\pi^* \alpha}{e_T(T_{I_\pi} \Hilb^n(\C^3))}\, [I_\pi]
\]
in $H^*_T(\Hilb^n(\C^3); \Q) \otimes_{H^*_T(\pt)} \mathrm{Frac}(H^*_T(\pt))$
(Atiyah--Bott~$1984$, \textit{Topology} $23$).
\end{remark}

\begin{definition}[Framed pre-$E_3$ structure on $\Hilb^\bullet(\C^3)$]
\label{def:prefactorization-hilb-c3}
\ClaimStatusDefinitional
The collection $\Hilb^\bullet(\C^3) = \bigsqcup_n \Hilb^n(\C^3)$ of
moduli of ideal subschemes carries a \emph{configuration-space
refinement}: for $U \subset \C^3$ open, write
\[
 \Hilb^\bullet(U) \;=\; \bigsqcup_n \{I \in \Hilb^n(\C^3)
 : \mathrm{supp}(\C[x, y, z] / I) \subset U\}.
\]
The assignment $U \mapsto H^*_T(\Hilb^\bullet(U); \Q)$ extends to a
pre-factorisation algebra (in the sense of Costello--Gwilliam
\textit{Factorization algebras in quantum field theory}, volume~$1$,
Chapter~$3$) on $\C^3$ when one passes to the completed version that
tracks ideal supports in an open $U$ rather than as abstract moduli.
For disjoint opens $U_1, U_2 \subset \C^3$, the restriction-of-support
map
\[
 \Hilb^\bullet(U_1 \sqcup U_2) \;\xrightarrow{\sim}\; \Hilb^\bullet(U_1)
 \times \Hilb^\bullet(U_2)
\]
sending $I$ to $(I \cap \C[U_1], I \cap \C[U_2])$ is an isomorphism on
the level of support decomposition; taking equivariant cohomology
yields the pre-$E_3$ factorisation structure map
\[
 H^*_T(\Hilb^\bullet(U_1)) \otimes H^*_T(\Hilb^\bullet(U_2))
 \;\xrightarrow{\;m_{U_1, U_2}\;}\;
 H^*_T(\Hilb^\bullet(U_1 \cup U_2)).
\]
\end{definition}

\begin{remark}[Side-by-side with Costello--Francis--Gwilliam $2026$]
\label{rem:cfg-2026-side-by-side}
Costello--Francis--Gwilliam~$2026$ Section~$1.1$ defines an $E_3$-algebra
structure on the local observables of topological Chern--Simons at gauge
$\widehat{\fgl}_1$ (or more generally at gauge $\mathfrak g$) via the
configuration-space model $\Conf_n(\R^3)$ of $n$-point configurations
in three-dimensional Euclidean space. The sphere-bundle fibre of this
configuration space is homotopy-equivalent to $S^{3 \cdot n - 1}$
minus the discriminant, and the $E_3$-structure is encoded by the
operadic composition of the little-$3$-disks operad $D_3$.
Definition~\ref{def:prefactorization-hilb-c3} above attaches to each
open $U \subset \C^3$ the corresponding ideal moduli
$\Hilb^\bullet(U)$; under the real-analytic homeomorphism
$\C^3 \simeq \R^6$, the little-$3$-disks operad $D_3$ acts on
$H^*_T(\Hilb^\bullet(-); \Q)$ via the configuration-space presentation
of the CoHA product. The $E_3$-structure of CFG $2026$ and the
$E_3$-structure of the present construction are the \emph{same}
$D_3$-operad action evaluated on two different targets: CFG $2026$
targets free-field topological Chern--Simons observables at gauge
$\widehat{\fgl}_1$; the present construction targets
$H^*_T(\Hilb^\bullet(-))$. The CFG $2026$ and SV identifications are
two faces of the same $E_3$-operadic action on cohomology.
\end{remark}

\subsection*{The Schiffmann--Vasserot convolution product: push-pull on $\mathrm{Flag}$}

\begin{definition}[Nakajima flag correspondence]
\label{def:nakajima-flag-correspondence}
\ClaimStatusDefinitional
For $n, m \geq 0$, let
\[
 \mathrm{Flag}^{n, m}(\C^3)
 \;=\; \{(I_1, I_2) \in \Hilb^n(\C^3) \times \Hilb^{n + m}(\C^3)
 : I_1 \supset I_2,\; \dim_\C I_1 / I_2 = m\}.
\]
Equivalently, $\mathrm{Flag}^{n, m}$ parametrises pairs of nested ideals
$I_1 \supset I_2$ with colength $n$ and $n + m$ respectively. The two
projections
\[
 \mathrm{Flag}^{n, m} \;\xrightarrow{\;p_1\;}\; \Hilb^n(\C^3),
 \qquad
 \mathrm{Flag}^{n, m} \;\xrightarrow{\;p_2\;}\; \Hilb^{n + m}(\C^3),
\]
are projective and equivariant. The virtual dimension of
$\mathrm{Flag}^{n, m}$ equals the sum of the virtual dimensions of
$\Hilb^n$ and $\Hilb^m$ by the perfect obstruction theory of Nakajima
(Nakajima--Yoshioka $2003$, \textit{Invent. Math.} $155$, \S$3$).
\end{definition}

\begin{definition}[CoHA product via push-pull]
\label{def:coha-product-c3}
\ClaimStatusDefinitional
The Schiffmann--Vasserot CoHA product on
$H^*_T(\Hilb^\bullet(\C^3); \Q)$ is the push-pull
\[
 \star \colon H^*_T(\Hilb^n) \otimes H^*_T(\Hilb^m)
 \;\xrightarrow{\;p_1^* \boxtimes \mathrm{id}\;}
 H^*_T(\mathrm{Flag}^{n, m})
 \;\xrightarrow{\;(p_2)_*\;}
 H^*_T(\Hilb^{n + m}),
\]
where $p_1^*$ is the flat pullback along the first projection and
$(p_2)_*$ is the proper pushforward along the second, in
$T$-equivariant Borel--Moore homology (Chriss--Ginzburg
\textit{Representation Theory and Complex Geometry}, Birkh\"auser
$1997$, \S$2.7$). The product is associative (Schiffmann--Vasserot
$2013$ Theorem~$4.1$). Explicitly, on equivariant cohomology classes
of the monomial-ideal fixed-point basis, the product is the
\emph{shuffle product} with rational kernel
\[
 \phi(u, v) \;=\; \frac{(u - v + \epsilon_1)(u - v + \epsilon_2)
 (u - v + \epsilon_3)}{(u - v)(u - v + \epsilon_1 + \epsilon_2)
 (u - v + \epsilon_2 + \epsilon_3)},
\]
in the torus-localised limit. (Schiffmann--Vasserot~$2013$
\textit{Publ. Math. IHES} $118$ Theorem~$1.1$; the rational shuffle
kernel originates in Feigin--Odesskii shuffle presentations, adapted
by SV to the CY locus $\epsilon_1 + \epsilon_2 + \epsilon_3 = 0$.)
\end{definition}

\begin{remark}[Associativity class]
\label{rem:coha-associativity-class}
The associativity of $\star$ lives at the level of Borel--Moore
homology; at the chain level (complexes of mixed Hodge modules) the
associator is a homotopy controlled by the triple-flag correspondence
\[
 \mathrm{Flag}^{n, m, p}(\C^3)
 \;=\; \{I_1 \supset I_2 \supset I_3 : \dim I_1 / I_2 = m,\;
 \dim I_2 / I_3 = p\}
\]
with its three-cartesian-square compatibility (Chriss--Ginzburg
\S$2.7$). The degree-$0$ associator coincides at cohomology; the
$\infty$-categorical coherence is provided by the configuration-space
presentation (cf.~Remark~\ref{rem:cfg-2026-side-by-side}), upgraded to
the Kontsevich--Tamarkin $E_3$-formality transfer of the bracket of
degree $-2$ on $\HH^\bullet(\Perf(\C^3)) = \Q[x, y, z]$ regarded as a
shifted symplectic dg-algebra (Pantev--To\"en--Vaqui\'e--Vezzosi
$2013$, \textit{Publ. Math. IHES} $117$). The CFG $2026$ approach
encodes this associator through the singular chain model of $D_3$;
the SV presentation encodes it through the triple-flag correspondence;
the two are intertwined via the equivalent factorisation-algebra
structure on $\Hilb^\bullet(\C^3)$.
\end{remark}

\subsection*{The Schiffmann--Vasserot theorem}

\begin{theorem}[Schiffmann--Vasserot $2013$]
\label{thm:sv-c3-wave10}
\ClaimStatusProvedElsewhere
With the notation above, there is an isomorphism of associative
$\Q$-algebras
\[
 \bigl(H^*_T(\Hilb^\bullet(\C^3); \Q),\; \star\bigr)
 \;\xrightarrow{\;\sim\;}\;
 Y^+(\widehat{\fgl}_1)
\]
onto the positive half of the affine Yangian of $\widehat{\fgl}_1$,
with parameters $(\epsilon_1, \epsilon_2, \epsilon_3)$ specialising to
the CY constraint $\epsilon_1 + \epsilon_2 + \epsilon_3 = 0$. The
isomorphism sends the fundamental generator $1 \in H^*_T(\Hilb^1(\C^3))
= H^*_T(\C^3)$ to the affine-Yangian generator $e_0 \in
Y^+(\widehat{\fgl}_1)$, and the shuffle product $\star$ to the
Feigin--Odesskii shuffle product on the $\widehat{\fgl}_1$-lattice.
(Schiffmann--Vasserot~$2013$ \textit{Publ. Math. IHES} $118$
Theorem~$1.1$.)
\end{theorem}

\begin{remark}[Four sources, one theorem]
\label{rem:sv-four-sources}
Theorem~\ref{thm:sv-c3-wave10} is one theorem with four historically
distinct presentations: (1) the Jordan-quiver critical CoHA
$\CoHA(Q = \text{Jordan}, W = \mathrm{tr}(XYZ - XZY))$ of
Kontsevich--Soibelman~$2010$ (arXiv:$1006.2706$) $\S 7.2$; (2) the
$3$d-moduli space Nakajima variety of ADHM data extended to $\C^3$
(Schiffmann--Vasserot $2013$ \S$3$); (3) the Maulik--Okounkov
$R$-matrix stable envelope on $\Hilb^\bullet(\C^2) \times \C$ lifted
to $\C^3$ in the CY limit (Maulik--Okounkov~$2012$ arXiv:$1211.1287$);
(4) the CFG $2026$ topological Chern--Simons observables at gauge
$\widehat{\fgl}_1$ evaluated on $\R^3 \simeq \C^3$. Sources (1)--(3)
output the same shuffle algebra $Y^+(\widehat{\fgl}_1)$; source (4)
recasts it as $E_3$-algebra observables of a topological field theory.
The four sources cross-check each other; no new mathematics is
claimed, but the CFG $2026$ reformulation exposes the $E_3$-operadic
content implicit in the SV proof.
\end{remark}

\subsection*{Chain-level refinement: $\PhiFA_3(\Perf(\C^3)) = \CoHA(\C^3)$ as $E_3$ factorisation algebra}

\begin{theorem}[SV as $E_3$ factorisation algebra on $\C^3$]
\label{thm:sv-e3-factorization-wave10}
\ClaimStatusConditional\textup{ (conditional on the Kontsevich--Tamarkin
$E_3$-formality transfer of Proposition~\ref{prop:phifa-infty1-kernel}
and on the Omega-background equivariant-localisation identification of
the Gerstenhaber bracket of degree $-2$ on
$\HH^\bullet(\Perf(\C^3))$)}
The Stage-$1$ output $\PhiFA_3(\Perf(\C^3))$ on $\C^3$, as $E_3$
holomorphic factorisation algebra, is computed by the
equivariant-cohomology assignment
\[
 U \;\longmapsto\; H^*_T(\Hilb^\bullet(U); \Q),
 \qquad U \subset \C^3\ \text{open},
\]
and the factorisation product $m_{U_1, U_2}$ of
Definition~\ref{def:prefactorization-hilb-c3} is the
Schiffmann--Vasserot CoHA product of Definition~\ref{def:coha-product-c3}
in the limit of disjoint support. Equivalently:
\[
 \PhiFA_3(\Perf(\C^3)) \;\simeq\; \CoHA(\C^3) \;=\; Y^+(\widehat{\fgl}_1)
 \qquad\text{in}\qquad E_3\text{-}\HolFA(\C^3).
\]
The CY constraint $\epsilon_1 + \epsilon_2 + \epsilon_3 = 0$ is the
holomorphic twist of $E_3$-operad to $E_3$-holomorphic operad (Costello--Li
$2016$ arXiv:$1605.09656$, $2020$ arXiv:$2001.07888$).
\end{theorem}

\begin{proof}[Proof sketch]
Three steps. First, HKR (Hochschild--Kostant--Rosenberg) on the
polynomial ring $\C[x, y, z]$ identifies
\[
 \HH^\bullet(\Perf(\C^3)) \;\simeq\; \mathrm{PV}^\bullet(\C^3)
 \;=\; \C[x, y, z] \otimes \Lambda[\partial_x, \partial_y, \partial_z],
\]
the polyvector fields on $\C^3$, with its native Schouten--Nijenhuis
Gerstenhaber bracket of degree $-2$. Second, the Kontsevich
$E_3$-formality theorem (Kontsevich~$1999$ arXiv:math/$9904055$;
Tamarkin~$2007$ arXiv:math/$9803025$) upgrades this Gerstenhaber
bracket to an $E_3$-algebra structure on $\mathrm{PV}^\bullet(\C^3)$,
canonically up to a $\mathrm{GRT}_1(\Q)$-torsor of Drinfeld associator
choices. Third, Costello--Gwilliam--Li holomorphic locality (Costello--Gwilliam
\textit{Factorization algebras in quantum field theory} volume~$2$
$\S 4.10$; Costello--Li~$2016$) assembles the $E_3$-algebra into a
holomorphic factorisation algebra on $\C^3$, carrying the extra datum
of the volume form $dx \wedge dy \wedge dz$. In the Omega-background
$T$-equivariant limit, the resulting $E_3$-holomorphic factorisation
algebra evaluated at open $U \subset \C^3$ is
$H^*_T(\Hilb^\bullet(U); \Q)$ with the push-pull product, by the
Feigin--Odesskii shuffle presentation of
$Y^+(\widehat{\fgl}_1)$ matching the SV identification of
Theorem~\ref{thm:sv-c3-wave10}. The composite is
$\PhiFA_3(\Perf(\C^3)) = Y^+(\widehat{\fgl}_1)$.
\end{proof}

\begin{remark}[Why the CY constraint enters]
\label{rem:cy-constraint-entry}
The CY constraint $\epsilon_1 + \epsilon_2 + \epsilon_3 = 0$ enters
twice. First, it is the condition that $\C^3$ is Calabi--Yau in the
$T$-equivariant sense: the canonical bundle $K_{\C^3} = \cO_{\C^3}$ is
$T$-equivariantly trivial precisely when the weights sum to zero.
Second, it is the condition that the Gerstenhaber bracket of degree
$-2$ on $\mathrm{PV}^\bullet(\C^3)$ is the Schouten--Nijenhuis bracket
rather than a twisted variant: for non-CY $\C^3$ equipped with a
shifted symplectic form of non-standard shift, the bracket lives in a
different Gerstenhaber operad. The two-parameter Omega-background
$(\epsilon_1, \epsilon_2)$ with $\epsilon_3 = -\epsilon_1 - \epsilon_2$
parametrises the deformation family of the CoHA away from the pure-CY
point; Schiffmann--Vasserot's rational shuffle kernel retains both
parameters. The chiral-algebra face (the evaluation module
$\cW_{1+\infty}[\lambda]$ for $\lambda \in \Q$) depends on $(\epsilon_1,
\epsilon_2)$ through the triality parameter $\lambda = \epsilon_1 /
\epsilon_2 + \epsilon_2 / \epsilon_1 - 2$ (Gaiotto--Rap\v{c}\'ak $2017$
arXiv:$1703.00982$; Proch\'azka--Rap\v{c}\'ak $2018$ arXiv:$1807.11304$).
\end{remark}

\subsection*{The Drinfeld double and the evaluation chain}

\begin{theorem}[Drinfeld double of the SV positive half]
\label{thm:drinfeld-double-sv-wave10}
\ClaimStatusProvedElsewhere
The positive half $Y^+(\widehat{\fgl}_1)$ admits a compatible negative
half $Y^-(\widehat{\fgl}_1) = Y^+(\widehat{\fgl}_1)^{\mathrm{op}}$ and
Cartan $Y^0(\widehat{\fgl}_1) = \Q[h_0, h_1, h_2, \ldots]$ (the mode
algebra of the $\widehat{\fgl}_1$ Heisenberg current). The Drinfeld
double
\[
 D\bigl(Y^+(\widehat{\fgl}_1)\bigr)
 \;=\; Y^+(\widehat{\fgl}_1) \otimes Y^0(\widehat{\fgl}_1)
 \otimes Y^-(\widehat{\fgl}_1)
 \;\simeq\; Y(\widehat{\fgl}_1)
\]
is the full affine Yangian of $\widehat{\fgl}_1$, a quasi-triangular
quasi-Hopf $\Q$-algebra. (Drinfeld $1988$ \textit{Dokl. Akad. Nauk SSSR}
$296$; Schiffmann--Vasserot $2013$ Theorem~$5.1$; Tsymbaliuk $2017$
arXiv:$1404.5240$ gives the explicit RTT-presentation.)
\end{theorem}

\begin{remark}[Three associativity classes]
\label{rem:three-associativity-classes-wave10}
Along the evaluation chain $\CoHA(\C^3) \hookrightarrow Y \to
\mathrm{End}(\cW_{1+\infty}[\lambda]\text{-vac})$ the associativity class
changes at each arrow. (A) $\CoHA(\C^3)$ is associative as a
graded $\Q$-algebra; its $E_3$-lift in
Theorem~\ref{thm:sv-e3-factorization-wave10} carries a homotopy
associator controlled by Kontsevich--Tamarkin $E_3$-formality. (B)
$Y^+ \hookrightarrow D(Y^+) = Y(\widehat{\fgl}_1)$ is the Drinfeld
double embedding; the target is quasi-Hopf, with braiding
$R$-matrix carrying the Yang--Baxter equation; the associativity
class upgrades from pure associative to braided quasi-Hopf. (C)
$\mathrm{ev}_\lambda \colon Y(\widehat{\fgl}_1) \to \mathrm{End}(\cW_{1+\infty}[\lambda]\text{-vac})$
is the evaluation into the mode algebra of a specific vertex algebra
module; the target $\cW_{1+\infty}[\lambda]$ is a vertex operator
algebra, carrying OPE associativity with locality (the Jacobi
identity). The three associativity classes sit in the sequence
\[
 \text{assoc.\ graded alg.}
 \;\hookrightarrow\; \text{quasi-Hopf braided alg.}
 \;\xrightarrow{\mathrm{ev}}\;
 \text{VOA mode algebra}.
\]
None of the three arrows is an isomorphism of the full structure; each
adds new data. In particular $\CoHA(\C^3) \neq \cW_{1+\infty}$: a graded
associative algebra is not a vertex algebra, and the triality parameter
$\lambda$ that labels the target module is absent from the CoHA side.
\end{remark}

\subsection*{The Maulik--Okounkov $R$-matrix on $\Hilb^\bullet(\C^3)$}

\begin{theorem}[Stable envelope on $\Hilb^\bullet(\C^3)$]
\label{thm:mo-stable-envelope-c3-wave10}
\ClaimStatusProvedElsewhere
For $\Hilb^\bullet(\C^3) = \bigsqcup_n \Hilb^n(\C^3)$, the
Maulik--Okounkov stable envelope $\mathrm{Stab}_\pm \colon
H^*_T(\Hilb^{\bullet, T}) \to H^*_T(\Hilb^\bullet)$ (with respect to
a generic chamber $\pm$ in the cocharacter lattice of $T$) produces
an $R$-matrix
\[
 R^{\mathrm{MO}}(u) \;=\; \mathrm{Stab}_+^{-1} \circ \mathrm{Stab}_-
 \;\in\; \mathrm{End}\bigl(H^*_T(\Hilb^\bullet(\C^3)) \otimes
 H^*_T(\Hilb^\bullet(\C^3))\bigr) [\![u^{-1}]\!],
\]
satisfying the Yang--Baxter equation. (Maulik--Okounkov~$2012$
arXiv:$1211.1287$; for the $\Hilb^\bullet(\C^2)$ case, which is the
reduction of the present $\Hilb^\bullet(\C^3)$ along the third
coordinate direction.)
\end{theorem}

\begin{remark}[MO $R$-matrix as residue of the UV gluing cocycle]
\label{rem:mo-r-matrix-residue-wave10}
The Vol III universal positive-geometry grammar (cache) realises
\[
 R^{\mathrm{MO}}(u) \;=\; \mathrm{Res}_{u = u_*}\,
 \phi^+_{\mathrm{UV}}(u),
\]
with $\phi^+_{\mathrm{UV}}$ the UV gluing cocycle of the positive half
$Y^+(X)$ across the equivariant chamber wall at $u = u_*$. For
$X = \C^3$ the cocycle is the rational shuffle kernel $\phi(u, v)$ of
Definition~\ref{def:coha-product-c3}, and the residue at $u = v$ is
precisely the MO $R$-matrix at zero spectral parameter.  The Yang--Baxter
axiom for $R^{\mathrm{MO}}(u)$ is the cocycle condition for
$\phi^+_{\mathrm{UV}}$ on triple intersections of chamber walls.
\end{remark}

\subsection*{The framing of CFG $2026$: topological Chern--Simons at $\widehat{\fgl}_1$}

\begin{theorem}[CFG $2026$ topological Chern--Simons and $\CoHA(\C^3)$]
\label{thm:cfg-2026-tcs-coha-wave10}
\ClaimStatusProvedElsewhere\textup{ (in CFG $2026$; stated here for the
side-by-side)}
Topological Chern--Simons theory on $\R^3$ at gauge group
$\widehat{\fgl}_1$ (or more generally at gauge $\mathfrak g$) produces
an $E_3$-algebra of local observables. Evaluated at the holomorphic
twist of $\C^3$ with gauge $\widehat{\fgl}_1$ and with volume form $dx \wedge dy \wedge dz$,
the observables are canonically identified with $\CoHA(\C^3) =
Y^+(\widehat{\fgl}_1)$; the $E_3$-algebra structure is the CoHA
associativity lifted by Kontsevich--Tamarkin $E_3$-formality, and the
holomorphic twist is the CY constraint $\epsilon_1 + \epsilon_2 +
\epsilon_3 = 0$.
\end{theorem}

\begin{remark}[The identification at the operadic level]
\label{rem:cfg-2026-operadic-identification-wave10}
At the level of operads, CFG $2026$'s $E_3$-algebra and the
$\PhiFA_3$-Stage-$1$ output on $\Perf(\C^3)$ are the same object:
both are the evaluation of the little-$3$-disks operad $D_3$ on the
polyvector-field Gerstenhaber algebra of $\C^3$ with degree-$(-2)$
Schouten bracket. The CFG $2026$ derivation uses the perturbative
expansion of topological CS with classical BV action
$S = \mathrm{tr}(A dA + \tfrac{2}{3} A^3)$; the SV derivation uses the
convolution on $\mathrm{Flag}^{n, m}(\C^3)$. The two derivations
produce the same output because both factor through the same
Kontsevich--Tamarkin associator on $\mathrm{PV}^\bullet(\C^3)$. CFG $2026$
is thus a \emph{parallel} source for the $E_3$-factorisation-algebra
structure on $\CoHA(\C^3)$, not an alternative: its primary technical
content is the Drinfeld associator choice compatible with
Costello--Li perturbative hCS, which coincides with the SV choice up
to $\mathrm{GRT}_1(\Q)$-torsor translation.
\end{remark}

\subsection*{Omega-background equivariant chain-level derivation}

\begin{definition}[Omega-background parameters]
\label{def:omega-background-wave10}
\ClaimStatusDefinitional
The \emph{Omega-background} on $\C^3$ is the pair $(\epsilon_1, \epsilon_2)$
of formal parameters, with the CY-imposed third parameter
$\epsilon_3 = -\epsilon_1 - \epsilon_2$. Physically, the background is
the Nekrasov $2004$ (arXiv:hep-th/$0306238$) $\Omega$-deformation of
$\mathcal N = 2$ gauge theory on $\C^2 \times S^1$ extended to the
CY$_3$ setting via the third torus direction.
$T$-equivariant cohomology of $\Hilb^\bullet(\C^3)$ localises to monomial
ideals (plane partitions), with weights in the two-parameter family.
\end{definition}

\begin{theorem}[Equivariant chain-level CoHA $=$ SV positive half]
\label{thm:equivariant-chain-coha-wave10}
\ClaimStatusProvedElsewhere
The $T$-equivariant Borel--Moore homology $H^*_T(\Hilb^\bullet(\C^3); \Q)$
with convolution product is canonically isomorphic, as a two-parameter
$\Q[\epsilon_1, \epsilon_2]$-algebra, to the positive half of the
affine Yangian of $\widehat{\fgl}_1$:
\[
 \bigl(H^*_T(\Hilb^\bullet(\C^3); \Q),\; \star\bigr)_{\epsilon_1, \epsilon_2}
 \;\xrightarrow{\;\sim\;}\; Y^+_{\epsilon_1, \epsilon_2}(\widehat{\fgl}_1).
\]
(Schiffmann--Vasserot $2013$ \textit{Publ. Math. IHES} $118$
Theorem~$1.1$; Nakajima $2014$ arXiv:$1401.6782$ for the $\C^2$ base
case with cyclic $T$-action; reduction to $\C^3$ via the third torus
direction.)
\end{theorem}

\begin{remark}[Two-parameter vs CY locus]
\label{rem:two-parameter-vs-cy-wave10}
The SV theorem operates with two free parameters $(\epsilon_1,
\epsilon_2)$ and recovers the CY locus as the hyperplane
$\epsilon_3 = -\epsilon_1 - \epsilon_2$. The full two-parameter family
is the Maulik--Okounkov / Nekrasov $\mathrm{U}(1)$ equivariant
$K$-theory setting; the CY locus is the restriction to the
Costello--Li holomorphic twist. The $\PhiFA_3$-Stage-$1$ output
depends only on the CY locus (the holomorphic-twist direction fixes
the third coordinate), but the full SV theorem is a statement about
the two-parameter family; $\Phi^{FA}_3(\Perf(\C^3))$ recovers the SV
output at the single CY point.
\end{remark}

\subsection*{The positive-geometry grammar of $Y^+(\C^3)$}

\begin{definition}[Universal positive-geometry presentation]
\label{def:universal-positive-geometry-c3-wave10}
\ClaimStatusDefinitional
Following the Vol III cache entry on the universal positive-geometry
grammar, we present $Y^+(\C^3)$ as
\[
 Y^+(\C^3) \;=\; H^*_{\mathrm{eq}}\bigl(\cM^+_{\mathrm{eff}}(\C^3),\;
 \phi_W\bigr)
\]
where $\cM^+_{\mathrm{eff}}(\C^3)$ is the positive cone of effective
classes in the moduli stack of quiver representations of the Jordan
quiver $Q$ with cubic potential $W = \mathrm{tr}(XYZ - XZY)$, and
$\phi_W$ is the vanishing-cycle sheaf of $W$ (Kontsevich--Soibelman
$2011$ arXiv:$1006.2706$ $\S 7.2$; Davison $2017$ arXiv:$1512.04179$
for the integrality refinement). The positive cone is the submonoid
generated by the fundamental effective class $[\C^3] \in
K_0(\Perf(\C^3))$; its $n$-fold product lands in $\Hilb^n(\C^3)$.
\end{definition}

\begin{remark}[Four equivariance strata locate $\C^3$]
\label{rem:four-equivariance-strata-c3-wave10}
Among the four equivariance strata of the Vol III cache — (i) toric
$T^d$; (ii) reduced $\C^\times + \mathrm{Aut}(X)$; (iii) orbifold
inertia $I(X/G)$; (iv) lattice-polarised period domain — the target
$\C^3$ sits in stratum (i), with full $T^3 = (\C^\times)^3$ equivariance
imposed by the ambient toric structure. The CY constraint reduces the
effective torus to $T^2 = \ker(\epsilon_1 + \epsilon_2 + \epsilon_3)$.
Stratum (i) is the complete stratum for SV and CFG $2026$: no orbifold
or lattice-polarisation structure needs to be added.
\end{remark}

\subsection*{The K\"ahler cone, BPS generating function, and PT $=$ DT on $\C^3$}

\begin{theorem}[PT $=$ DT $=$ MacMahon on $\C^3$]
\label{thm:macmahon-c3-wave10}
\ClaimStatusProvedElsewhere
The generating function of $T$-equivariant Donaldson--Thomas invariants
of $\C^3$, evaluated at the point-count partition function, is the
MacMahon function
\[
 Z_{\mathrm{DT}}^{\mathrm{eq}}(\C^3)(q) \;=\; \prod_{n \geq 1} (1 - q^n)^{-n}
 \;=\; M(q).
\]
By MNOP--$\mathrm{I}$ (arXiv:math/$0312059$) Theorem $1$ this coincides
with the PT (Pandharipande--Thomas $2007$ arXiv:$0707.2348$) invariant
generating function. By Schiffmann--Vasserot $2013$ Theorem~$1.1$ this
equals the graded dimension of $Y^+(\widehat{\fgl}_1)$:
\[
 \sum_n \dim_\Q \CoHA(\C^3)_n\, q^n \;=\; M(q).
\]
\end{theorem}

\begin{remark}[MacMahon as $E_3$ bar Euler characteristic]
\label{rem:macmahon-bar-euler-wave10}
Equivalently (Vol III toric\_cy3\_coha Section on $\C^3$ as MacMahon,
Proposition~\ref{prop:toric-kappa-table} of the target chapter), the
$E_1$-bar Euler characteristic of $A_{\C^3}$ also equals $M(q)$. This
is \emph{not} accidental: the graded dimension of $\CoHA(\C^3) =
Y^+(\widehat{\fgl}_1)$ and the $E_1$-bar Euler characteristic of the
SV-image $\cW_{1+\infty}$-vacuum generating function coincide because
both are the MacMahon partition function of plane partitions, which
Schiffmann--Vasserot identify as the common bilateral graded structure
via the Nakajima basis of $H^*_T(\Hilb^\bullet(\C^3))$ indexed by
plane partitions.  The identification is $\kappa_{\mathrm{ch}}(\C^3) =
3/2$ in the CY convention, read off the cubic-potential graded Euler
characteristic (stratum (i) of the cache).
\end{remark}

\subsection*{Five ghost theorems and their beating mathematical cores}

\begin{remark}[Ghost theorem 1: $\CoHA(\C^3) = \cW_{1+\infty}$]
\label{rem:ghost-theorem-1-wave10}
The \emph{ghost theorem}: $\CoHA(\C^3) = \cW_{1+\infty}$, as associative
$\Q$-algebras. \emph{Beating core}: the correct identification is
$\CoHA(\C^3) = Y^+(\widehat{\fgl}_1) \hookrightarrow Y(\widehat{\fgl}_1)
\xrightarrow{\mathrm{ev}_\lambda} \mathrm{End}(\cW_{1+\infty}[\lambda]
\text{-vacuum})$, and $\cW_{1+\infty}$ is the evaluation image on a
specific vacuum module, not an isomorphic copy of the CoHA. The
vertex-algebra face $\cW_{1+\infty}$ is the \emph{representation
target}; it is a vertex operator algebra with OPE associativity,
locality, and triality parameter $\lambda = \epsilon_1 / \epsilon_2 +
\epsilon_2 / \epsilon_1 - 2$; the CoHA is associative-graded, with
triality parameter absent. The three associativity classes of
Remark~\ref{rem:three-associativity-classes-wave10} prevent the
conflation.
\end{remark}

\begin{remark}[Ghost theorem 2: $\CoHA(\C^3) = Y(\widehat{\fgl}_1)$]
\label{rem:ghost-theorem-2-wave10}
\emph{Ghost theorem}: $\CoHA(\C^3) = Y(\widehat{\fgl}_1)$ full Yangian.
\emph{Beating core}: CoHA is only the positive half; the full Yangian
is the Drinfeld double $D(Y^+(\widehat{\fgl}_1))$. The negative half
$Y^- = \CoHA(\C^3)^{\mathrm{op}}$ is the opposite CoHA; the Cartan
$Y^0$ is generated by the $\widehat{\fgl}_1$ modes $\{h_n\}_{n \in \Z}$.
The positive half alone carries only the positive-frequency
$\widehat{\fgl}_1$ modes $h_0, h_1, h_2, \ldots$; to get all modes one
must Drinfeld-double.
\end{remark}

\begin{remark}[Ghost theorem 3: PhiFA$_3$ of $\Perf(\C^3)$ gives a vertex algebra]
\label{rem:ghost-theorem-3-wave10}
\emph{Ghost theorem}: $\PhiFA_3(\Perf(\C^3))$ is already a vertex
algebra (a $2$d chiral algebra). \emph{Beating core}:
$\PhiFA_3(\Perf(\C^3))$ is a \emph{$3$-dimensional} $E_3$-holomorphic
factorisation algebra on $\C^3$, not a $2$d chiral algebra. A vertex
algebra emerges only after Stage-$2$ specialisation
$\mathrm{Sp}_{\Sigma_2, C}$ to a curve $C$, pushing along a codimension-one
cycle $\Sigma_2 \subset \C^3$; without a choice of $\Sigma_2$ there is
no vertex-algebra output. The $\C^3$ case has no canonical
compactification to a curve, so the specialisation is parametrised by
the choice of $\Sigma_2$ (the two-dimensional cycle in $\C^3$), and
different choices give different vertex algebras: the
$\Sigma_2 = \{z = 0\}$ choice gives the $\C^2 \subset \C^3$ reduction
with vertex algebra $\cW_{1+\infty}[\lambda]$ via $\mathrm{ev}_\lambda$
evaluation, with $\lambda$ determined by the two surviving parameters
$(\epsilon_1, \epsilon_2)$.
\end{remark}

\begin{remark}[Ghost theorem 4: six routes to $\CoHA(\C^3)$ are six $\PhiFA_3$-applications]
\label{rem:ghost-theorem-4-wave10}
\emph{Ghost theorem}: the Schiffmann--Vasserot, RSYZ, MNOP, MO-stable-envelope,
and CFG $2026$ derivations of $Y^+(\widehat{\fgl}_1)$ are all different
applications of $\PhiFA_3$. \emph{Beating core}: all are the \emph{same}
$\PhiFA_3(\Perf(\C^3)) = Y^+(\widehat{\fgl}_1)$, derived through
different operadic or moduli-theoretic presentations. They are five
\emph{consistency cross-checks} on a single mathematical object, not
five distinct $\Phi$-applications. The output is one $E_3$-holomorphic
factorisation algebra on $\C^3$; the input is one Calabi--Yau category
$\Perf(\C^3)$; the derivations differ in their technical route.
\end{remark}

\begin{remark}[Ghost theorem 5: CoHA is independent of the choice of associator]
\label{rem:ghost-theorem-5-wave10}
\emph{Ghost theorem}: the SV isomorphism $\CoHA(\C^3) =
Y^+(\widehat{\fgl}_1)$ is canonical, independent of the Drinfeld
associator. \emph{Beating core}: the $E_3$-lift of the associative
CoHA depends on a choice of Drinfeld associator up to
$\mathrm{GRT}_1(\Q)$-torsor (Kontsevich $1999$ formality; Tamarkin
$2007$ deformation quantisation; Willwacher $2011$ arXiv:$1009.1654$
on $\mathrm{GRT}_1$). The \emph{cohomological} (or
\emph{derived-category}) CoHA is canonical; the \emph{$E_3$-factorisation-algebra}
lift is canonical up to contractible choice, with the $\mathrm{GRT}_1(\Q)$-torsor
parametrising the choice of Drinfeld associator. CFG $2026$ and SV
use compatible associator choices; the choice is invisible at cohomology
level but shows up as a specific identification of the $E_3$-structure
with the Costello--Li holomorphic-twist data.
\end{remark}

\subsection*{Stable core after five attack-heal cycles}

\begin{proposition}[Stable core of the chain-level SV]
\label{prop:stable-core-wave10}
\ClaimStatusProvedHere
The following are unconditional:
\begin{enumerate}[label=\textup{(\roman*)}]
 \item $\CoHA(\C^3) = H^*_T(\Hilb^\bullet(\C^3); \Q)$ with push-pull
  product on $\mathrm{Flag}^{n, m}$ is an associative graded $\Q$-algebra,
  canonically isomorphic to $Y^+(\widehat{\fgl}_1)$ (Schiffmann--Vasserot
  $2013$ Theorem~$1.1$; unconditional).
 \item $Y^+(\widehat{\fgl}_1)$ sits inside $Y(\widehat{\fgl}_1) =
  D(Y^+(\widehat{\fgl}_1))$ as the Drinfeld-double positive half
  (Schiffmann--Vasserot $2013$ Theorem~$5.1$; unconditional).
 \item $Y(\widehat{\fgl}_1)$ evaluates onto $\mathrm{End}(\cW_{1+\infty}[\lambda]
  \text{-vacuum})$ via the Gaiotto--Rap\v{c}\'ak evaluation
  $\mathrm{ev}_\lambda$ (Gaiotto--Rap\v{c}\'ak $2017$; Proch\'azka--Rap\v{c}\'ak
  $2018$; unconditional).
 \item The graded dimension of $\CoHA(\C^3)$ equals the MacMahon function
  $M(q) = \prod_{n \geq 1}(1 - q^n)^{-n}$, which equals the
  $T$-equivariant DT partition function of $\C^3$ (MNOP--$\mathrm{I}$
  $2006$ Theorem~$1$; Schiffmann--Vasserot $2013$ Theorem~$1.2$;
  unconditional).
 \item The Maulik--Okounkov stable envelope on $\Hilb^\bullet(\C^3)$
  produces an $R$-matrix satisfying Yang--Baxter (Maulik--Okounkov $2012$
  Theorem~$1$; unconditional).
\end{enumerate}
The following are conditional on the $\PhiFA_d$-assembly
(Proposition~\ref{prop:phifa-infty1-kernel}, which is proved in the
infinity-categorical framework):
\begin{enumerate}[label=\textup{(\alph*)}]
 \item $\PhiFA_3(\Perf(\C^3)) \simeq \CoHA(\C^3) = Y^+(\widehat{\fgl}_1)$
  as $E_3$-holomorphic factorisation algebras on $\C^3$. Conditional on
  the Kontsevich--Tamarkin $E_3$-formality transfer of the degree-$(-2)$
  Gerstenhaber bracket on $\HH^\bullet(\Perf(\C^3))$ to an $E_3$-algebra
  on $\mathrm{PV}^\bullet(\C^3)$, and on the Costello--Gwilliam--Li
  assembly of this $E_3$-algebra into a holomorphic factorisation
  algebra on $\C^3$.
 \item The CFG $2026$ topological Chern--Simons $E_3$-algebra of local
  observables at gauge $\widehat{\fgl}_1$ agrees with $\CoHA(\C^3) =
  Y^+(\widehat{\fgl}_1)$ at operadic level. Conditional on the CFG
  $2026$ proof unrolling through the same $\mathrm{GRT}_1(\Q)$-torsor
  of associator choices used by Kontsevich--Tamarkin and compatible
  with Costello--Li perturbative hCS.
 \item The $R$-matrix $R^{\mathrm{MO}}(u)$ equals the residue of the
  UV gluing cocycle $\phi^+_{\mathrm{UV}}(u)$ of $Y^+(\C^3)$ across
  the equivariant chamber wall $u = u_*$. Conditional on the universal
  positive-geometry grammar of the Vol III cache, which is stated at
  stratum (i) for $\C^3$ and holds as a theorem in this stratum.
\end{enumerate}
\end{proposition}

\subsection*{Pathways forward}

\begin{remark}[Three concrete paths forward]
\label{rem:three-paths-forward-wave10}
\emph{Path $1$}: extend the CFG $2026$ identification to the full toric
CY$_3$ family by the RSYZ analogue of Theorem~\ref{thm:cfg-2026-tcs-coha-wave10}.
For a toric CY$_3$ $X_\Sigma$ without compact $4$-cycles, topological
Chern--Simons at gauge $\widehat{\fg}_{Q_X}$ (the affine super Yangian
attached to the toric quiver $Q_X$) produces an $E_3$-algebra of local
observables that should agree with $\CoHA(Q_X, W_X)$. This requires
extending Costello--Li perturbative hCS from $\mathrm{U}(1)$ gauge to
affine-Yangian gauge; open.

\emph{Path $2$}: promote the $E_3$-factorisation algebra structure of
Theorem~\ref{thm:sv-e3-factorization-wave10} to a stable $\infty$-category
of CoHA modules. The target is a stable $\infty$-category
$\mathrm{Mod}_{\PhiFA_3(\Perf(\C^3))}$ equipped with an $E_3$-monoidal
structure; applications to the Drinfeld-centre computation and to
Conjecture CY-C (existence of $G(\C^3) = Y(\widehat{\fgl}_1)$ as a
chiral quantum group) would follow.

\emph{Path $3$}: compute the \v{C}ech descent on the $\C^3$ factorisation
algebra along a resolution $\widetilde{\C^3} \to \C^3$ (for instance
$\mathrm{Tot}(\cO_{\bP^2}(-3)) \to \C^3 / \Z_3$ for local $\bP^2$). The
descent datum is a sheaf of $Y^+$-algebras on $\widetilde{\C^3}$ with
orbifold-inertia twist; the target chapter's Section on three-chart
McKay orbifold gluing (Theorem~\ref{thm:local-p2-mckay}) would be
refined to the $E_3$-factorisation level through this path.
\end{remark}

\subsection*{Side-by-side comparison: CFG $2026$ vs Schiffmann--Vasserot $2013$}

\begin{remark}[Side-by-side of the two derivations]
\label{rem:side-by-side-derivations-wave10}
\begin{center}
\begin{tabular}{l|l|l}
\textbf{Structure} & \textbf{CFG $2026$ derivation} & \textbf{Schiffmann--Vasserot $2013$} \\
\hline
Operad & $D_3 \simeq E_3$ from $\Conf_n(\R^3)$ & Convolution on $\mathrm{Flag}^{n, m}(\C^3)$ \\
Input & Top.\ CS action $\mathrm{tr}(A dA + \tfrac{2}{3} A^3)$ & Jordan quiver, $W = \mathrm{tr}(XYZ - XZY)$ \\
Target & Observables of top.\ CS at $\widehat{\fgl}_1$ & $H^*_T(\Hilb^\bullet(\C^3); \Q)$ \\
Parameter & Perturbative $\hbar$ expansion & Equivariant $(\epsilon_1, \epsilon_2, \epsilon_3)$ \\
Associator & Kontsevich configuration integral & Triple-flag compatibility \\
Output & $E_3$-algebra of observables & $\CoHA(\C^3) = Y^+(\widehat{\fgl}_1)$ \\
CY constraint & Top.\ twist $\to$ hol.\ twist & $\epsilon_1 + \epsilon_2 + \epsilon_3 = 0$ \\
Identification & Operadic via $D_3$ action & Shuffle algebra via Feigin--Odesskii \\
Cross-check & MO stable envelope $R$-matrix & MacMahon generating function \\
\hline
\end{tabular}
\end{center}
Both derivations factor through the same Kontsevich--Tamarkin
$\mathrm{GRT}_1(\Q)$-torsor of Drinfeld associator choices; both produce
the same mathematical output $Y^+(\widehat{\fgl}_1)$. CFG $2026$'s
contribution is making the $E_3$-operadic structure explicit at the
level of observables, previously implicit in the SV convolution
formula.
\end{remark}

\subsection*{Cycle 1 attack: does $\PhiFA_3$ really output the CoHA?}

\begin{remark}[Attack and heal, cycle 1]
\label{rem:cycle-1-attack-heal-wave10}
\emph{Attack}: The Stage-$1$ output $\PhiFA_3(\Perf(\C^3))$ is an
$E_3$-holomorphic factorisation algebra on $\C^3$ coming from
$\HH^\bullet(\Perf(\C^3))$ with its degree-$(-2)$ Gerstenhaber bracket.
The CoHA, by contrast, lives on $\Hilb^\bullet(\C^3)$, a moduli of
ideals, not on $\C^3$ itself. The putative identification
$\PhiFA_3(\Perf(\C^3)) \simeq \CoHA(\C^3)$ conflates two different
moduli: the derived category $\Perf(\C^3)$ and the derived moduli of
ideals $\Hilb^\bullet(\C^3)$. The identification cannot hold.

\emph{Surviving core}: the Hilbert scheme $\Hilb^\bullet(\C^3)$ is the
moduli of objects in $\Perf(\C^3)$ (more precisely, of ideal
subsheaves with given generic rank), and its equivariant cohomology
is a derived invariant of $\Perf(\C^3)$. The two moduli are not the
same space, but their cohomologies $\HH^\bullet(\Perf(\C^3))$ and
$H^*_T(\Hilb^\bullet(\C^3))$ are bilaterally related: by HKR the former
is the polyvector fields on $\C^3$; by Nakajima the latter is a
$\Z_{\geq 0}$-graded sum of cohomologies indexed by ideal colength $n$.
The $E_3$-factorisation-algebra structure descends from the
$\HH^\bullet$-Gerstenhaber bracket via Costello--Li assembly, and the
resulting factorisation algebra has $n$-point observables given by
$H^*_T(\Hilb^n)$ in the $T$-equivariant limit. The identification is
\emph{not} a direct equality of moduli but an identification of the
$E_3$-algebra of observables with the convolution CoHA on
$\Hilb^\bullet$.

\emph{Heal}: Theorem~\ref{thm:sv-e3-factorization-wave10} states the
identification at the level of $E_3$-holomorphic factorisation algebras,
not at the level of underlying moduli spaces. The composite
$\Phi^{FA}_3(\Perf(\C^3)) \simeq \CoHA(\C^3) = Y^+(\widehat{\fgl}_1)$
holds as an equivalence in $E_3\text{-}\HolFA(\C^3)$, with the RHS
realised by $U \mapsto H^*_T(\Hilb^\bullet(U); \Q)$ under the
support-refinement of Definition~\ref{def:prefactorization-hilb-c3}.
\end{remark}

\subsection*{Cycle 2 attack: the $\mathrm{GRT}_1(\Q)$-torsor}

\begin{remark}[Attack and heal, cycle 2]
\label{rem:cycle-2-attack-heal-wave10}
\emph{Attack}: the Kontsevich--Tamarkin $E_3$-formality depends on a
choice of Drinfeld associator in the $\mathrm{GRT}_1(\Q)$-torsor of
rational associators. Different associator choices give different
$E_3$-structures on $\mathrm{PV}^\bullet(\C^3)$; therefore the
identification $\PhiFA_3(\Perf(\C^3)) \simeq \CoHA(\C^3)$ cannot be
canonical — there is an associator ambiguity.

\emph{Surviving core}: the $\mathrm{GRT}_1(\Q)$-torsor acts on the
$E_3$-lift of a fixed Gerstenhaber algebra, but it does \emph{not}
change the cohomology: $H^*(\mathrm{PV}^\bullet(\C^3)) =
\mathrm{PV}^\bullet(\C^3)$ is associator-independent. The SV
identification $\CoHA(\C^3) = Y^+(\widehat{\fgl}_1)$ holds at the
level of the associative (Kontsevich) graded algebra, which is
associator-independent. The associator ambiguity only affects the
specific $E_3$-chain-level model; at cohomology (Borel--Moore homology
of $\Hilb^\bullet(\C^3)$) the identification is canonical.

\emph{Heal}: Theorem~\ref{thm:sv-e3-factorization-wave10} asserts the
equivalence in $E_3\text{-}\HolFA(\C^3)$ \emph{up to contractible
choice} (Kontsevich $1999$ formality theorem: the space of
$E_3$-quasi-isomorphism classes on $\mathrm{PV}^\bullet(\C^3)$ is
contractible modulo $\mathrm{GRT}_1(\Q)$-torsor action). CFG $2026$
and SV agree on this contractible space, up to the
$\mathrm{GRT}_1(\Q)$-torsor translation which is the distance between
their respective Drinfeld associator choices (the CFG $2026$
perturbative CS associator vs the SV configuration-integral
associator). The torsor is a $\Q$-vector space of dimension bounded
by the Brown--Schnetz mixed-Tate $\zeta$-values in weight $\leq 4$;
neither derivation pins down a specific element, but both land in the
same orbit.
\end{remark}

\subsection*{Cycle 3 attack: is the Omega-background mandatory?}

\begin{remark}[Attack and heal, cycle 3]
\label{rem:cycle-3-attack-heal-wave10}
\emph{Attack}: SV requires $T$-equivariance with Omega-background
parameters $(\epsilon_1, \epsilon_2)$; without equivariance,
$H^*(\Hilb^\bullet(\C^3); \Q)$ is not the affine Yangian
$Y^+(\widehat{\fgl}_1)$. But $\PhiFA_3(\Perf(\C^3))$ is defined in
the non-equivariant setting — the Gerstenhaber bracket on
$\HH^\bullet(\Perf(\C^3))$ does not mention the torus $T$. Therefore
the identification
$\PhiFA_3(\Perf(\C^3)) \simeq \CoHA(\C^3)$ must fail non-equivariantly.

\emph{Surviving core}: the Calabi--Yau structure $\epsilon_1 +
\epsilon_2 + \epsilon_3 = 0$ is not an equivariance datum; it is the
holomorphic-twist constraint that picks out CY$_3$ from the
two-parameter $\mathrm{U}(1)$-equivariant deformation. Without the
$T$-action, the CoHA is the non-equivariant Borel--Moore homology
$H^{\mathrm{BM}}_*(\Hilb^\bullet(\C^3); \Q)$, which at infinite
dimension is a completion / localisation issue but \emph{not} a
failure of the identification. The $T$-equivariance provides a
convenient basis (the plane-partition basis of
Remark~\ref{rem:hilb-c3-torus-action}), and the two-parameter family
is the deformation; the identification holds at all parameter values
simultaneously, with the non-equivariant limit being the specialisation
$\epsilon_1 = \epsilon_2 = 0$.

\emph{Heal}: the $\Omega$-background is a technical convenience, not
a mandatory ingredient. $\PhiFA_3(\Perf(\C^3))$ and $\CoHA(\C^3)$
both carry the two-parameter family as a deformation; the SV
identification is a statement about the entire deformation family,
and the CY locus $\epsilon_1 + \epsilon_2 + \epsilon_3 = 0$ is a
hyperplane inside it. The non-equivariant limit is a point on this
hyperplane, and the identification restricts to it cleanly.
\end{remark}

\subsection*{Cycle 4 attack: the Stage-$2$ specialisation}

\begin{remark}[Attack and heal, cycle 4]
\label{rem:cycle-4-attack-heal-wave10}
\emph{Attack}: the target of $\Phi_3$ is a chiral algebra on a curve;
the output on $\C^3$ is an $E_3$-holomorphic factorisation algebra.
These are different categories of output. The SV identification
$\CoHA(\C^3) = Y^+(\widehat{\fgl}_1)$ gives a graded associative
algebra, not a chiral algebra; therefore the CoHA cannot be the
$\Phi_3$-output. The chain-level evaluation chain runs through
$\cW_{1+\infty}[\lambda]$, a vertex algebra; the CoHA is not the
VOA. The claim $\PhiFA_3(\Perf(\C^3)) = \CoHA(\C^3)$ contradicts the
expected chiral-algebra output.

\emph{Surviving core}: the two-stage factorisation $\Phi_3 =
\mathrm{Sp}_{\Sigma_2, C} \circ \PhiFA_3$ has \emph{Stage $1$}
$\PhiFA_3$ outputting an $E_3$-factorisation algebra on $\C^3$
(not a chiral algebra on a curve) and \emph{Stage $2$}
$\mathrm{Sp}_{\Sigma_2, C}$ specialising to a curve $C$ via pushforward
over a codimension-$2$ cycle $\Sigma_2$. The CoHA is the Stage-$1$
output, not the full $\Phi_3$-output. The vertex-algebra output
$\cW_{1+\infty}[\lambda]$ is the Stage-$2$ specialisation to
$C = \mathbb P^1$ with $\Sigma_2 = \{z = 0\} \subset \C^3$ and the
$\lambda$-parameter reading the residual two-parameter data.

\emph{Heal}: the SV identification takes place at Stage~$1$:
$\PhiFA_3(\Perf(\C^3)) = \CoHA(\C^3) = Y^+(\widehat{\fgl}_1)$. This is
not a vertex algebra; it is an $E_3$-factorisation algebra. The
vertex algebra emerges only after Stage~$2$ specialisation to a curve,
which is the evaluation chain
$Y^+ \hookrightarrow Y \xrightarrow{\mathrm{ev}_\lambda}
\mathrm{End}(\cW_{1+\infty}[\lambda]\text{-vacuum})$. The two stages
are kept separate, following the Vol III cache discipline on two-stage
factorisation and Pattern $273$.
\end{remark}

\subsection*{Cycle 5 attack: the CFG $2026$ is a separate mathematical programme}

\begin{remark}[Attack and heal, cycle 5]
\label{rem:cycle-5-attack-heal-wave10}
\emph{Attack}: CFG $2026$ treats topological $3$d Chern--Simons as
a \emph{topological} field theory and produces an $E_3$-algebra of
observables by topological operadic methods. SV treats the convolution
on $\Hilb^\bullet(\C^3)$ as an \emph{algebraic} construction using
equivariant cohomology. These are different mathematical settings;
the agreement is a coincidence.

\emph{Surviving core}: both CFG $2026$ and SV compute the same
underlying $E_3$-algebra of observables, namely the $E_3$-lift of the
polyvector-field Gerstenhaber algebra on $\C^3$. The CFG $2026$
topological framework uses the little-$3$-disks operad $D_3$ on the
topological target $\R^3 \simeq \C^3$; SV uses the convolution-style
product on $\mathrm{Flag}^{n, m}$ in the algebraic category. Both are
presentations of the same $E_3$-algebra object because both factor
through the Kontsevich--Tamarkin transfer $\mathrm{PV}^\bullet(\C^3)
\to E_3(\mathrm{PV}^\bullet(\C^3))$.

\emph{Heal}: CFG $2026$ is a parallel presentation, not an independent
construction. The agreement is a consequence of the uniqueness
(up to $\mathrm{GRT}_1(\Q)$-torsor) of the $E_3$-lift of the Gerstenhaber
bracket on $\mathrm{PV}^\bullet(\C^3)$. CFG $2026$ does mathematical
work that SV does not: it exposes the perturbative CS-action origin of
the associator choice and connects it to Costello's holomorphic-twist
programme. The combined CFG $2026$ + SV picture is stronger than either
alone.
\end{remark}

\subsection*{A final remark on the cache evaluation chain}

\begin{remark}[Evaluation chain, chain-level version]
\label{rem:evaluation-chain-wave10-final}
The cache evaluation chain
\[
 \CoHA(\C^3) \;=\; Y^+(\widehat{\fgl}_1)
 \;\hookrightarrow\; Y(\widehat{\fgl}_1)
 \;\xrightarrow{\;\mathrm{ev}_\lambda\;}\;
 \mathrm{End}(\cW_{1+\infty}[\lambda]\text{-vacuum})
\]
lifts to the chain level as follows. (A) $\CoHA(\C^3) = Y^+(\widehat{\fgl}_1)$
is an $E_3$-algebra in $\mathrm{Ch}(\Q)$ with the Kontsevich--Tamarkin
associator; its associative reduction is Schiffmann--Vasserot's
shuffle algebra. (B) $Y^+ \hookrightarrow Y$ is the Drinfeld-double
embedding; both are associative $\Q$-algebras at the cohomology level,
with the $E_3$-lift of $Y^+$ extending to a quasi-triangular
quasi-Hopf structure on $Y$ with braiding the Yang--Baxter $R^{\mathrm{MO}}(u)$.
(C) $\mathrm{ev}_\lambda$ is the evaluation into the mode algebra
of $\cW_{1+\infty}[\lambda]$, realising the Yangian image as the
positive-frequency and negative-frequency modes of the chiral
currents; this is a homomorphism of associative algebras with vertex-algebra
image, not a morphism of $E_3$-algebras (vertex algebras sit in
$E_1$-holomorphic factorisation algebras on curves). The three steps
thus involve three category changes: $E_3$-algebra on $\C^3$ $\to$
quasi-Hopf algebra (Drinfeld double) $\to$ vertex algebra on $\bP^1$;
each step preserves some structure and forgets others.

The Vol III cache observation that ``$\cW_{1+\infty}$ is the evaluation
image, not a completion or deformation'' is the statement that the
arrow $\mathrm{ev}_\lambda$ is a homomorphism, not an injection of
$\mathrm{End}$s. The positive-modes of $\cW_{1+\infty}[\lambda]$ are
the image of $Y^+$; the negative-modes are the image of $Y^-$; the
intersection on the vacuum vector is trivial (the Cartan acts by
central-charge scalars). The vertex-algebra structure emerges from
the OPE of the mode algebra, not from the Yangian structure.
\end{remark}

\subsection*{Summary diagram}

\begin{remark}[Master diagram, chain level]
\label{rem:master-diagram-wave10}
\[
 \begin{array}{ccccc}
 \Perf(\C^3) & \xrightarrow{\HH^\bullet, \text{HKR}} & \mathrm{PV}^\bullet(\C^3) & \xrightarrow{\mathrm{KT}} & E_3(\mathrm{PV}^\bullet(\C^3)) \\
 & & & & \downarrow \mathrm{Costello-Li} \\
 & & & & E_3\text{-}\HolFA(\C^3) \\
 & & & & \|\;\mathrm{Thm.}~\ref{thm:sv-e3-factorization-wave10} \\
 & & & & \CoHA(\C^3) \;=\; Y^+(\widehat{\fgl}_1) \\
 & & & & \downarrow D\;\text{(Drinfeld double)} \\
 & & & & Y(\widehat{\fgl}_1) \\
 & & & & \downarrow \mathrm{ev}_\lambda\;\text{(Stage-$2$ spec.)} \\
 & & & & \mathrm{End}(\cW_{1+\infty}[\lambda]\text{-vac}) \\
 & & & & \downarrow\;\text{VOA extraction} \\
 & & & & \cW_{1+\infty}[\lambda] \\
 \end{array}
\]
The vertical pathway from $E_3\text{-}\HolFA(\C^3)$ down to
$\cW_{1+\infty}[\lambda]$ is the composite of: (i) the SV identification
(one equality); (ii) the Drinfeld-double embedding; (iii) the
Gaiotto--Rap\v{c}\'ak evaluation; (iv) the VOA mode-algebra
extraction. Only (i) is an equivalence in $E_3\text{-}\HolFA(\C^3)$;
(ii) is a quasi-Hopf embedding; (iii) is a VOA evaluation; (iv) is a
VOA mode decomposition. The chiral-algebra object $\cW_{1+\infty}[\lambda]$
emerging at the bottom is the Stage-$2$ specialisation output of
$\Phi_3(\Perf(\C^3))$ at the specific choice of cycle
$\Sigma_2 = \{z = 0\} \subset \C^3$ and curve $C = \bP^1$; different
Stage-$2$ choices on the same Stage-$1$ output give different vertex
algebras, each labelled by a triality parameter $\lambda$.
\end{remark}

\subsection*{Cross-reference to the target chapter}

\begin{remark}[Chapter integration]
\label{rem:chapter-integration-wave10}
The chain-level refinement above lifts the existing
Theorem~\ref{thm:sv-c3} of \texttt{toric\_cy3\_coha.tex} from a
purely associative statement (SV identification of graded associative
algebras) to an $E_3$-factorisation-algebra statement
(Theorem~\ref{thm:sv-e3-factorization-wave10}), integrated with the
Stage-$1$ output $\PhiFA_3(\Perf(\C^3))$ of the two-stage factorisation
of Proposition~\ref{prop:phifa-infty1-kernel}. The existing evaluation
chain of Remark~\ref{rem:evaluation-chain-c3} is extended to the chain
level via Remark~\ref{rem:evaluation-chain-wave10-final}, and the
five ghost theorems of this note place guards against the conflations
identified by the Vol III cache (AP-CY entries on CoHA $\neq$ vertex
algebra, $Y^+ \neq Y$, and six routes to $G(X)$ as different
constructions rather than different $\Phi$-applications).

The chapter already contains the CFG $2026$ reference in the
introduction; the present wave file makes the CFG $2026$ / SV /
Kontsevich--Tamarkin / Costello--Li four-way agreement explicit as a
statement about $\mathrm{GRT}_1(\Q)$-torsor of Drinfeld associators
and places the four derivations as parallel sources of the single
$E_3$-algebra object $Y^+(\widehat{\fgl}_1)$.
\end{remark}
